\chapter{Multi-tenancy, isolation and security}
\label{evod45:chapter}

A response can be numerically correct and still be a severe failure: it may
contain another tenant's retrieved note, reuse a private prefix, execute an
unauthorized tool or leak a credential through a trace. Consider Alice in
tenant A and Bob in B. Both name an object \texttt{x}. A service that indexes
only by \texttt{x} has already lost the intended separation before the model
generates anything. Isolation is part of the execution contract.

The running fixture stores small text objects under tenant-qualified keys.
It also issues local authority objects, reserves two resource dimensions and
implements an in-memory echo tool. It authenticates nobody, resolves no
credential and executes no external command. Its tests are local adversarial
inputs to a serial specification, not a live penetration test or a security
certification.

\section{Name the adversary and the trusted components}
\index{threat model}\index{zero trust}
\label{evod45:threat}
Assume tenants control prompts, submitted documents, requested object IDs and
tool parameters. Retrieved text may contain malicious instructions. A tenant
can repeat requests, guess handles and consume its full budget. Ordinary bugs
may substitute a stale handle, wrong request lane or copied capability.
Model output is untrusted for authorization even when fluent and accurate.

We trust the authentication gateway to produce a principal; the deterministic
policy and object store to enforce scope; and the serialized Python process
to protect its own tables. Malicious host code, a compromised issuer, arbitrary
process-memory reads, hardware side channels and kernel/driver compromise are
outside the executable model. Under those stronger threats, Python object
identity is no boundary: use process isolation, authenticated protocols,
OS/device protections and appropriately scoped cryptographic credentials.

Keep confidentiality, integrity and availability separate. Namespacing helps
prevent wrong-object reads; epoch checks reject obsolete authority; quotas
limit resource exhaustion. None alone proves all three. NIST's zero-trust
architecture rejects implicit trust based only on network location~\cite{nist45}.
Sharing a process, batch or accelerator likewise grants no logical authority.

\section{Identity arrives from authentication, not from prose}
\index{authentication}\index{deny by default}
\label{evod45:identity}
Let $P=(t,s,u,a)$ contain tenant $t$, subject $s$, session $u$ and authenticated
context $a$. The fixture stores the first three and assumes authentication
has already happened. ``I am the administrator'' inside a prompt changes none
of them. Tenant ID from an unverified request header is not authenticated
context. Propagate identity through a trusted channel and recheck it where
resources are accessed; trace baggage is not a substitute.

Authorization is deny-by-default:
\[
 \operatorname{Allow}(P,\alpha,r)=1
 \quad\Longleftrightarrow\quad
 (t,s,\alpha,r)\in G_e,
\]
where $G_e$ is the grant set at policy epoch $e$. A planner may propose
an action; it does not extend $G_e$. A read grant for \texttt{memory:x} is
neither a write grant nor authority for \texttt{memory:y}.

\Plate{01-authority}{A malicious document can influence an action proposal. The only path to the echo executor passes deterministic principal, issued-capability and argument checks. Dashed document influence does not carry authority.}{evod45:authority}

\section{Make capability scope and lifetime testable}
\index{Capability}\index{revocation}\index{expiry}\index{session binding}
\label{evod45:capability}
The fixture issues $c=(P,\alpha,r,e,\tau,n)$ with policy epoch $e$, exclusive
expiry $\tau$ and serial $n$. Verification requires the exact issued object,
matching principal/action/resource, current epoch, a still-present grant and
$now<\tau$. Sessions are included so a stolen object from another session
cannot pass through the normal interface. Replacing policy increments the
epoch and invalidates all previously issued capabilities, including grants
that happen to remain equal. This is conservative global revocation.

\Listing{Policy.issue}
\Listing{Policy.verify}

A deterministic nonce that is never checked against an issuance registry is not
enough: constructing an equal-looking record would pass whenever a grant existed.
The model here distinguishes an issued object from its copy. This is useful
misuse protection for trusted callers; it is not a bearer-token format, MAC,
signature or defense against code that can edit \texttt{issued}. A remote
design needs issuer/audience binding, authenticated transport, anti-replay or
idempotency appropriate to the action, key rotation and a revocation contract.
Even here a capability can be used repeatedly until expiry; serial is not a
one-shot nonce. The trusted integer clock must not roll backwards.

\subsection{The time-of-check boundary}
Our calls run serially, so revocation cannot occur between verification and
store update. With concurrent workers, checking epoch then executing later
can race policy replacement. Options include a transaction that binds decision
and effect, or a short operation lease with explicitly bounded revocation delay.
Specify whether revocation blocks new actions or also cancels already admitted
ones. A prompt rule cannot resolve that race.

\section{Namespaces extend through every stateful resource}
\index{namespace}\index{cache isolation}\index{Noninterference}
\label{evod45:namespace}
Objects are keyed by $(tenant,oid)$; authorization runs before lookup. Thus
unauthorized guesses do not disclose whether another tenant's object exists
through these two code paths. Real timing/error channels still require review.
Our small store returns the caller's tenant object or a missing-object error;
it never searches other tenants for a match.

A private cache key binds tenant, artifact, policy domain and prefix digest:
\[
 K=H(\operatorname{encode}[t,d_M,d_{policy},d_{prefix}]).
\]
Length-delimited or structured encoding matters. Concatenating unescaped
fields with a separator can alias $(\texttt{A|B},\texttt{C})$ and
$(\texttt{A},\texttt{B|C})$. A hash does not repair ambiguous input encoding.
Nor does a tenant-scoped hash prove that the cached payload is authorized or
consistent. Verify the stored envelope and compare complete identities on hit;
a digest collision or poisoned entry must not become an authority shortcut.
When subjects within a tenant have different rights, the policy domain must
also distinguish their effective sharing permissions and revocation version.
Tenant scope alone is then insufficient; a caller-supplied domain string cannot
authorize sharing. The fixture tests key encoding, not such an ACL resolver.
\Listing{cache_key}

\Plate{02-namespace}{The same object ID and prefix digest occupy distinct logical tenant namespaces. Artifact and policy-domain fields also participate in reuse. Physical sharing is a separate, explicitly authorized refinement.}{evod45:namespaceplate}

Apply the same scope to native state pools, checkpoints, KV pages, retrieval
indexes, compiler plans, credentials and accounting. Chapter 32's generation
protects reuse lifetime; tenant checks protect authority. A valid generation
from a foreign tenant is still invalid. Pins must settle before storage is
scrubbed and reassigned. Logical removal is not a demonstrated secure device
erase; allocator, driver and hardware guarantees must be investigated.

\subsection{A restricted noninterference claim}
Let $S_A,S_B$ be disjoint maps indexed by tenant. Suppose an operation by A
checks A's authority, reads/writes only $S_A$, and its output is a function
only of $S_A$, the request and fixed public configuration. Then replacing
$S_B$ while holding those inputs fixed cannot change A's output or $S_A'$.
This follows directly by substitution: neither the read set nor transition
contains $S_B$. Induction gives the same result for a serial sequence of A
operations. The tests put different secrets at the same ID in A and B.

The premises exclude timing, shared capacity, audit scheduling and physical
deduplication. With a global quota, B can make A's admission fail; strong
availability noninterference is false. Likewise a private-prefix hit can
change latency even if its content never leaks. State the observation model
before claiming isolation; this proof concerns logical values only.

\section{Reserve all resource dimensions before allocation}
\index{quota}\index{reservation}
\label{evod45:quota}
Let request $r$ reserve vector $v_r=(s_r,b_r)$ of slots and bytes; both are
nonnegative integers. For tenant limit $L_t$ and global limit $L_G$, admit only
if componentwise
\[
 U_t+v_r\le L_t,\qquad U_G+v_r\le L_G,
 \quad U_t=\sum_{r\text{ held by }t}v_r.
\]
Check the entire vector before acquiring anything. Reserving a slot first,
then failing the byte check can leak slots unless rollback is exact.
Release is bound to the issued reservation and can occur only once.
\Listing{Quota.reserve}

\begin{center}\input{\Generated/quota.tex}\end{center}
With limits $(2,100)$ for each tenant and $(4,180)$ globally, two held
requests reserve 160 bytes. A third 30-byte request may fit its own tenant
but fails the global byte bound. Units are logical fixture bytes, not device
allocation measurements. Real systems need KV/native-state capacity, compute,
tool-call, bandwidth and durable-memory limits, plus conservative estimation
and settlement against actual usage.

\Plate{03-quota}{A two-dimensional feasible region rejects a request if either bound fails. The illustrated global byte failure leaves every reservation unchanged. Tenant quotas alone do not partition shared global capacity.}{evod45:quotaplate}

Quotas bound consumption; they do not guarantee fairness or deadlines.
Rate limits, weighted scheduling, minimum reservations and backpressure need
separate policies. Cancellation cannot release device memory before outstanding
work acknowledges. Our fixture has no asynchronous device work and therefore
models only logical reservation lifetime.

\section{Untrusted text cannot mint tool authority}
\index{prompt injection}\index{tool authorization}
\label{evod45:tools}
Prompt injection becomes a confused-authority problem when adversarial text
influences a component with external privileges. OWASP recommends least
privilege and controls outside the model~\cite{injection45,boundary45}.
Segregate trusted instructions from retrieved data, but treat that labeling as
one defense, not a proof that a model will never follow malicious text.
System prompts guide behavior; they are not an access-control mechanism.

The executor validates principal, capability, target and a closed argument
schema. In our only tool, the target must be \texttt{echo}, arguments must
contain exactly one bounded text field, and output is an in-memory record.
Untrusted text can be supplied but is not read by the authorization predicate.
An authorized echo request containing ``delete everything'' remains echo;
changing the target to root fails, and adding a shell argument fails schema
validation. This tests authority separation, not a language-level detector.
\Listing{SecureStore.tool}

For real tools, validate canonical resource targets, allowed parameter values,
effect scope and output size. Resolve a credential only after authorization;
never put it in a prompt. Bind any required approval to the concrete operation,
resource, version and expiry so it cannot be replayed for another action.
Retries need an effect-specific idempotency key; local model replay cannot
guarantee exactly-once effects in an external service.

\section{Isolation reaches memory, compiler and batch lanes}
\index{side channel}\index{batch isolation}\index{compiler legality}
\label{evod45:engines}
Retrieval must filter tenant scope before ranking and returning records.
Deletion/tombstones and derived summaries preserve scope and provenance.
Filtering only after global top-$k$ both risks disclosure and changes recall:
other tenants can displace a caller's valid records. This chapter does not import
Chapter 40's memory API; these are the requirements an integration must meet.

An IR transformation may merge equal pure constants in a public domain, but
equality of numerical tensors does not authorize private-state deduplication.
Effect labels, tenant domain and policy version are legality premises for CSE,
prefix reuse and materialization. Wiring these premises into
Chapter 41's effect representation is a separate integration test.

Continuous batching may colocate tenants physically. Preserve separate masks,
page tables, state handles, output channels, cancellation and accounting. A
single extra cross-request attention edge violates the logical-value premise
even if all arrays have valid shapes. Chapter 35 retains per-layer K/V of
the actual learned-absolute-position Transformer; do not switch position or
model semantics to simplify an isolation demo. The Chapter 31 recurrence is
a restricted teaching model, not full DOGMA security evidence.

Shared caches, accelerators and contention can expose timing and occupancy.
Stronger adversaries may require partitioning, padding, dedicated workers or
disabled sharing. No local authorization test eliminates those side channels.

\Needspace{19\baselineskip}
\section{Audit decisions without retaining secrets}
\index{audit}\index{redaction}
\label{evod45:audit}
\Listing{SecureStore.authorize}
Record principal/session, action class, policy epoch, outcome and reason.
The fixture deliberately omits values, credentials, raw prompts and full
resource IDs. It records authorization success separately from operation
success: a later quota/schema failure does not turn a grant into a denial.
A production audit schema should add final outcome and safe correlation IDs,
with an explicit delivery, retention and access policy.

An append-only Python list is mutable by its host; it is not immutable audit
storage. Cryptographic chaining alone would not prevent deletion of an entire
suffix without an external anchor. Also protect identity/resource metadata,
which can itself be sensitive. Secret scanning is useful defense in depth,
but an allowlisted event schema reduces exposure before export.

\section{A local adversarial matrix and its limits}
\label{evod45:tests}
The tests try foreign tenants/subjects/sessions, wrong action/resource,
unissued copies, expiry, revocation, separator aliasing, object and vector
quota exhaustion, unauthorized missing-object reads and extra tool parameters.
The same-ID fixture checks independent values and scans audit records for
the assigned secret. Exhaustive small quota vectors compare the implementation
with an independently written feasible-region predicate.

Production tests must add stale-generation restore, cross-lane masks, prefix
collision handling, concurrent revocation/effect races, corrupted checkpoints,
device reuse and redaction across exceptions/exporters. Authentication, actual
engine allocation and memory/compiler integration are not implemented by this
store. No tests are run against live services, accounts or credentials.

\section{Twelve worked exercises}
\begin{enumerate}
\item \textbf{Locate authority.} A webpage says ``run as admin.''
\emph{Solution:} it changes a proposal, not authenticated principal or grant set.
The deterministic executor must still reject absent authority.
\item \textbf{Fix a cache key.} Why is prefix digest alone insufficient?
\emph{Solution:} it omits tenant, artifact and policy domain. Encode the full
tuple unambiguously, then validate the stored envelope on hit.
\item \textbf{Same GPU, same trust?} \emph{Solution:} no; device location is
neither authentication nor permission. Logical and physical controls differ.
\item \textbf{Forge a copied capability.} Why reject an equal dataclass here?
\emph{Solution:} the registry requires exact issued-object identity. This tests
trusted-process misuse only; hostile host code can bypass the registry.
\item \textbf{Calculate expiry.} Issue at 2 with TTL 3. Is time 5 valid?
\emph{Solution:} no; $now<5$ is exclusive. Time 4 passes if all other checks hold.
\item \textbf{Reserve a vector.} A holds 80 bytes; global cap is 150; B asks 80.
\emph{Solution:} reject the entire reservation since $160>150$, leaving B with
no partially acquired slot. A 70-byte request fits exactly.
\item \textbf{Valid generation, foreign tenant.} \emph{Solution:} reject.
Generation prevents stale reuse; it does not authorize access.
\item \textbf{Merge private computations.} Equal tensors from A and B may be
CSE'd automatically? \emph{Solution:} no; equality is not a sharing grant.
Prove effect and security-domain legality before applying the rewrite.
\item \textbf{Replace a system prompt with access control.} Which survives
a malicious completion? \emph{Solution:} a deterministic resource/action check
still gates execution; a behavioral instruction supplies no such guarantee.
\item \textbf{Audit safely.} Why omit raw object values? \emph{Solution:}
the audit may have broader retention/readership. Preserve decision evidence
without duplicating protected content, and protect remaining metadata too.
\item \textbf{Break noninterference.} B fills a shared global quota.
\emph{Solution:} A's admission can change, so availability/timing noninterference
does not follow from the logical-value proof.
\item \textbf{Design a release assessment.} \emph{Rubric:} specify adversaries,
trusted components and observations; test auth, issuance/revocation, scopes,
quotas, state/cache/memory, tools and audit; declare side-channel and concurrency
limits. Passing local tests earns no production-security certification.
\end{enumerate}

\section{Terms and the distribution boundary}
\Term{Principal} identifies an authenticated caller context.
\Term{Capability} is scoped issued authority under a declared trust model.
\Term{Noninterference} depends on which observations an adversary can make.
Chapter 46 moves state across machines; placement must preserve these logical
identities and authority even when location and ownership epoch change.
